\begin{abstract}
Modern provenance-based Intrusion Detection Systems (IDS) for detecting Advanced Persistent Threats (APTs) suffer from a critical deployment bottleneck: they require curated, pristine benign training data to build normal profiles. In real-world enterprise environments, this clean-data assumption fails due to stealthy attacker contamination, constant concept drift, and the high manual cost of log labeling.
We present the first provenance-based IDS framework that infers latent cyber risk directly from uncurated, unlabeled streaming logs without any clean-baseline collection period.
Our approach projects raw telemetry into a 12-dimensional Behavioral State Space using online Structural Causal Models (SCMs) learned from potentially contaminated streaming provenance graphs. Within this state space, system dynamics are evaluated in parallel by three reasoning operators: (1) Local State Reasoning ($R_L$) to measure instantaneous mechanism violations, (2) Temporal State Reasoning ($R_T$) via sequence modeling to capture state-transition risk, and (3) Behavioral State Reasoning ($R_B$) via a Behavioral Pattern Reasoner (BPR) to classify complex attack manifolds. While the state representation is learned in a strictly unsupervised manner from raw streams, the BPR and risk fusion layers are calibrated using a small validation set with sparse threat labels, yielding a hybrid semi-supervised inference pipeline with split conformal prediction for distribution-free false-alarm control. The key novel components introduced are: (i) a formal 12-dimensional causal Behavioral State Space grounded in five invariant coordinate groups derived from online SCM outputs; (ii) a Causal Intervention Score (CIS) rooted in graph boundary theory to distinguish localized noise from multi-hop attack propagation; and (iii) a multi-perspective parallel reasoning architecture that resolves the model-mismatch failures of pure causal detectors.

We evaluate on five diverse benchmarks (OpTC, BETH, StreamSpot, TC3, and NODLINK) against standard baselines (Isolation Forest, LOF, One-Class SVM, and Autoencoders). Our framework achieves near-perfect threat inference on simulated benchmarks (AUC 0.9837 on TC3 and 0.9784 on NODLINK) and Linux host telemetry (AUC 0.9901 on BETH). On real enterprise Windows logs (OpTC), we achieve an AUC of 0.9628 (F1 32.9\% at 5\% FPR), surpassing the original pure causal baseline by +12.7\% and exceeding both Isolation Forest (0.5968) and autoencoder (0.9364) baselines. The framework improves StreamSpot AUC from 0.4991 to 0.6909, confirming that behavioral state reasoning provides a robust, drift-adaptive, and deployable framework for streaming enterprise environments.
\end{abstract}

\begin{IEEEkeywords}
Intrusion Detection, Advanced Persistent Threats, Causal Discovery, Behavioral State Space, Stacked Ensemble, Conformal Prediction, Provenance Graphs
\end{IEEEkeywords}
